% 第3章 环境（Environments）
\chapter{环境}

上一章我们看到：把稀土离子当作完全自由的离子——既不彼此相互作用、也不与周围环境
相互作用——就能推断许多含稀土晶体的磁性质。然而对某些晶体中的磁性离子，通常不能
忽略这类相互作用；对许多材料，它们很大且重要。本章考虑原子与其近邻环境之间的
相互作用；下一章讨论晶体中磁性原子与相邻磁性原子之间的直接磁相互作用。

\section{晶体场}

要理解晶体造成的局域环境对原子能级的影响，首先需要回顾原子轨道的形状。s、p、d
轨道电子密度的角分布示于图 3.1。该图只画出各轨道波函数的角度部分，还有径向部分
（详见附录 C）。只有 s 轨道是球对称的，其余轨道都有显著的角依赖。这一点至关重要，
因为局域环境往往不是球对称的，不同的轨道于是会有不同的表现。

\begin{figure}
\centering
\includegraphics[width=0.8\textwidth]{fig3_01.pdf}
\caption{s、p、d 轨道的角分布。$\mathrm{d}_{z^2}$ 与 $\mathrm{d}_{x^2-y^2}$
能级合称 $e_g$ 能级；$\mathrm{d}_{xy}$、$\mathrm{d}_{xz}$、$\mathrm{d}_{yz}$
能级合称 $t_{2g}$ 能级。$\mathrm{d}_{z^2}$ 轨道有时也记作
$\mathrm{d}_{3z^2-r^2}$。}
\end{figure}

\subsection{晶体场的起源}

晶体场（crystal field）是由晶体中相邻原子产生的电场。在晶体场理论中，相邻轨道
被模型化为负的点电荷；对这一近似的改进是配体场理论（ligand field theory）——它
本质上是分子轨道理论的推广，关注中心离子的 d 轨道及其与周围离子（配体，ligands）
轨道的交叠。晶体场效应的大小与性质关键取决于局域环境的对称性。一个常见的情形是
八面体环境：在许多过渡金属化合物中，过渡金属离子位于八面体中心，每个顶点上是一个
氧之类的离子。此时晶体场主要来自氧轨道中带负电的电子的静电排斥。八面体与四面体
环境的示意图见图 3.2。

\begin{figure}
\centering
\includegraphics[width=0.5\textwidth]{fig3_02.pdf}
\caption{处于（a）八面体与（b）四面体环境中的金属原子 M。八面体环境常见于许多过渡
金属氧化物：氧阴离子位于八面体顶点，金属原子居中。四面体环境可方便地取立方体的
相间顶点来描述（如图）。}
\end{figure}

d 轨道分为两类：指向 $x$、$y$、$z$ 轴之间的 $t_{2g}$ 轨道（即 $\mathrm{d}_{xy}$、
$\mathrm{d}_{xz}$、$\mathrm{d}_{yz}$）和沿轴指向的 $e_g$ 轨道（沿 $z$ 轴有瓣的
$\mathrm{d}_{z^{2}}$ 轨道，以及沿 $x$、$y$ 轴都有瓣的 $\mathrm{d}_{x^{2}-y^{2}}$
轨道）。设想把一个含十个 d 电子的阳离子放到半径为 $r$ 的均匀带负电球的中心。
在这个球对称环境中五个 d 轨道全部简并（当然电荷的存在抬高了整个系统的能量）。
现在设想球上的电荷聚集成六个分立的点电荷，各位于一个正八面体的顶点、仍在球面上。
所有 d 轨道的总电子能量不变，但 d 轨道不再简并——这就造出了八面体环境。

为了说明环境对不同轨道的影响不同，看图 3.3：它以平面图（图 3.2(a) 在 $xy$ 面上的
投影）画出八面体环境中的两个 d 轨道。晶体场主要由相邻原子的 p 轨道产生。显然
$\mathrm{d}_{xy}$ 轨道（图 3.3(a)）与相邻 p 轨道的交叠比 $\mathrm{d}_{x^{2}-y^{2}}$
（图 3.3(b)）小，因而静电能量更低。

\begin{figure}
\centering
\includegraphics[width=0.45\textwidth]{fig3_03.pdf}
\caption{晶体场源于静电相互作用。八面体环境中 $\mathrm{d}_{xy}$ 轨道（a）的能量比
$\mathrm{d}_{x^{2}-y^{2}}$（b）低。}
\end{figure}

在八面体环境中，相邻正电荷聚于 $(\pm r,0,0)$、$(0,\pm r,0)$、$(0,0,\pm r)$，指向
$x$、$y$、$z$ 轴之间的三个轨道 $\mathrm{d}_{xy}$、$\mathrm{d}_{yz}$、$\mathrm{d}_{zx}$
能量降低，而沿 $x$、$y$、$z$ 轴指向的 $\mathrm{d}_{z^2}$ 与 $\mathrm{d}_{x^2-y^2}$
能量升高。五个能级于是按图 3.4(a) 劈裂：三重 $t_{2g}$ 能级下降，二重 $e_g$ 能级
上升。

\begin{figure}
\centering
\includegraphics[width=0.5\textwidth]{fig3_04.pdf}
\caption{（a）八面体与（b）四面体环境中的晶体场。}
\end{figure}

若局域环境不是八面体对称，晶体场甚至可能反向作用。例如四面体环境中，沿轴指向的
轨道恰好最大程度避开了位于（描写四面体的）立方体四个顶点上的原子的电荷密度
（见图 3.2(b)），故此时二重 $e_g$ 能级更低（图 3.4(b)）。

\begin{figure}
\centering
\includegraphics[width=0.45\textwidth]{fig3_05.pdf}
\caption{$3d^{6}$ 离子（如 $\mathrm{Fe^{2+}}$）（a）高自旋（弱场）与（b）低自旋
（强场）情形的电子组态。}
\end{figure}

处理 3d 电子不满的过渡金属离子时，已有的电子先填最低的（此时即 $t_{2g}$）能级再填
$e_g$ 能级。但填充的确切次序取决于晶体场能量与“把两个电子放进同一轨道”的库仑
代价之间的竞争——后者称为成对能（pairing energy），通常为正。若晶体场能量小于
成对能（弱场情形），电子加入时先逐个占据各轨道，然后才有轨道被双占据；反之若
晶体场能量大于成对能（强场情形），电子会先双重占据低能轨道，然后再去攀登高能
轨道的“险峰”。这些情形示于图 3.5。

八面体环境中，填入 1、2、3、8、9、10 个电子时方式毫无歧义；有趣的是 4、5、6、7
个电子的情形。下面以六电子为例（对应 $\mathrm{Fe^{2+}}$ 离子，见图 3.5）。

\begin{example}{3.1}
$\mathrm{Fe^{2+}}$ 离子有 $3d^{6}$ 壳层，成对能为正。弱场情形：每个轨道先填一个，
还剩一个电子，它只好不情愿地与某个 $t_{2g}$ 电子配对。结果四个未成对电子、
$S=2$（图 3.5(a)），称为高自旋组态。强场情形：六个电子全被塞进三个 $t_{2g}$
轨道，$e_g$ 轨道空着；没有未成对电子，$S=0$（图 3.5(b)），称为低自旋组态。
\end{example}

在某些含 $\mathrm{Fe^{2+}}$ 的材料中，可以用温度、压强甚至光照在低自旋与高自旋
组态之间引发自旋转变。

\subsection{轨道淬灭}

3d 离子的预期磁基态列于表 3.1。按 2.5 节的洪特规则计算 S、L、J 是轻而易举的事。
但如该节所述，预言的磁矩 $g_J[J(J+1)]^{1/2}$ 并不总与实验一致——例外是 $3d^{5}$
与 $3d^{10}$：它们的电子壳层半满或全满，$L=0$。

\begin{table}
\centering
\caption{用洪特规则得到的 3d 离子磁基态。对每个离子列出壳层组态与预言的基态
S、L、J，并给出按洪特规则计算的 $p=\mu_{\mathrm{eff}}/\mu_{\mathrm{B}}$。此值记作
$p_1=g_J[J(J+1)]^{1/2}$；下一列是由含相应离子的顺磁盐测量得到的实验值 $p_{\mathrm{exp}}$。
它与 $p_2=2[S(S+1)]^{1/2}$ 符合好得多——后者假定轨道淬灭，即 L=0、J=S、$g_J=2$。}
\small
\begin{tabular}{lcccccccc}
\toprule
离子 & 壳层 & $S$ & $L$ & $J$ & 谱项 & $p_1$ & $p_{\mathrm{exp}}$ & $p_2$\\
\midrule
$\mathrm{Ti^{3+}},\mathrm{V^{4+}}$ & $3d^{1}$ & $\frac{1}{2}$ & 2 & $\frac{3}{2}$ & ${}^{2}D_{3/2}$ & 1.55 & 1.70 & 1.73\\
$\mathrm{V^{3+}}$ & $3d^{2}$ & 1 & 3 & 2 & ${}^{3}F_{2}$ & 1.63 & 2.61 & 2.83\\
$\mathrm{Cr^{3+}},\mathrm{V^{2+}}$ & $3d^{3}$ & $\frac{3}{2}$ & 3 & $\frac{3}{2}$ & ${}^{4}F_{3/2}$ & 0.77 & 3.85 & 3.87\\
$\mathrm{Mn^{3+}},\mathrm{Cr^{2+}}$ & $3d^{4}$ & 2 & 2 & 0 & ${}^{5}D_{0}$ & 0 & 4.82 & 4.90\\
$\mathrm{Fe^{3+}},\mathrm{Mn^{2+}}$ & $3d^{5}$ & $\frac{5}{2}$ & 0 & $\frac{5}{2}$ & ${}^{6}S_{5/2}$ & 5.92 & 5.82 & 5.92\\
$\mathrm{Fe^{2+}}$ & $3d^{6}$ & 2 & 2 & 4 & ${}^{5}D_{4}$ & 6.70 & 5.36 & 4.90\\
$\mathrm{Co^{2+}}$ & $3d^{7}$ & $\frac{3}{2}$ & 3 & $\frac{9}{2}$ & ${}^{4}F_{9/2}$ & 6.63 & 4.90 & 3.87\\
$\mathrm{Ni^{2+}}$ & $3d^{8}$ & 1 & 3 & 4 & ${}^{3}F_{4}$ & 5.59 & 3.12 & 2.83\\
$\mathrm{Cu^{2+}}$ & $3d^{9}$ & $\frac{1}{2}$ & 2 & $\frac{5}{2}$ & ${}^{2}D_{5/2}$ & 3.55 & 1.83 & 1.73\\
$\mathrm{Zn^{2+}}$ & $3d^{10}$ & 0 & 0 & 0 & ${}^{1}S_{0}$ & 0 & 0 & 0\\
\bottomrule
\end{tabular}
\end{table}

偏差的原因是：对 3d 离子，晶体场相互作用远强于自旋--轨道相互作用。于是 2.5 节
表述的洪特第三条规则——它基于“自旋--轨道相互作用是仅次于静电效应的最重要能量项”
——实际上是错的。数据显示这些体系转而选择 $L=0$（从而 $J=S$、$g_J=2$）的基态，
即
\begin{equation*}
\mu_{\mathrm{eff}} = 2\mu_{\mathrm{B}}\sqrt{S(S+1)}.\tag{3.1}
\end{equation*}
如表 3.1 所示，这一式子与实验符合得好得多。这个效应称为轨道淬灭（orbital
quenching），此时称轨道矩被淬灭（quenched）。对 4f 离子，轨道在空间上离核较近、
藏在 5s 与 5p 壳层之下，晶体场项的重要性小得多，洪特第三条规则得以遵守。更高过渡金属
离子（4d、5d 系列）的情形不那么明确：较重的离子自旋--轨道劈裂更大，晶体场与自旋--
轨道相互作用的效应可能不相上下。

对自旋--轨道相互作用可以忽略的 3d 离子，让我们更细致地考察轨道淬灭。八面体环境
中的晶体场等于常数加上正比于
$x^{4}+y^{4}+z^{4}-\frac{3}{5}r^{4}+O(r^{6}/a^{6})$ 的项（见习题 3.2），因此是实
函数。晶体场哈密顿量的典型特点是可以用实函数（不含微分算符）表达，故其本征函数
全为实函数。而总角动量算符 $\hat{\mathbf{L}}$ 是厄米的，本征值为实，但算符本身是
纯虚的。\fn{1}{式 C.23 表明角动量算符 $\hat{\mathbf{L}}=-\mathrm{i}\hbar\,\mathbf{r}\times\nabla$ 是纯虚的。}于是，晶体场劈裂所实现的非简并基态 $|0\rangle$（非简并——这样才
没法与另一个基态做线性组合来“耍花招”）必须是实函数；因此 $\langle 0|\hat{\mathbf{L}}|0\rangle$
因 $\hat{\mathbf{L}}$ 纯虚而为纯虚数；但 $\hat{\mathbf{L}}$ 厄米，$\langle 0|\hat{\mathbf{L}}|0\rangle$
又必须是纯实数。既纯实又纯虚，只能是
\begin{equation*}
\langle 0|\hat{\mathbf{L}}|0\rangle = 0\tag{3.2}
\end{equation*}
即非简并态的轨道角动量各分量全部被淬灭。

这个结果还可换个角度看。$\hat{L}_z$ 的本征值为 $m_l$ 的轨道态有方位角依赖
$e^{\mathrm{i}m_l\phi}$。本征函数须为实，这要求把 $\pm m_l$ 的态做线性组合，\fn{2}{这是为了构造方位角依赖如 $\cos(m_l\phi)$ 或 $\sin(m_l\phi)$ 的态。注意 $\cos(m_l\phi)=\frac{1}{2}(e^{\mathrm{i}m_l\phi}+e^{-\mathrm{i}m_l\phi})$，$\sin(m_l\phi)=\frac{1}{2\mathrm{i}}(e^{\mathrm{i}m_l\phi}-e^{-\mathrm{i}m_l\phi})$。}所得态的 $\langle L_z\rangle=0$。轨道淬灭的一种半经典解读是：轨道角动量在晶体场中进动，
大小不变而各分量平均为零。

实际上轨道角动量可能没有被完全淬灭，因为即便在 3d 离子中自旋--轨道相互作用也非
全然可忽略。它可以作为微扰加入，混入角动量非零的态，得到淬灭的基态但 $g$ 因子
不完全等于纯自旋值 2——与 2 的差别反映混入的 $L>0$ 态。这还会使 $g$ 因子略有
各向异性：其值随磁场相对晶体轴的方向略有不同。

\begin{figure}
\centering
\includegraphics[width=0.42\textwidth]{fig3_06.pdf}
\caption{$\mathrm{Mn^{3+}}$（$3d^{4}$）的扬--特勒效应。八面体配合物（左）可以畸变
（右），使 $t_{2g}$ 与 $e_g$ 能级劈裂。畸变之所以降低能量，是因为单占据的 $e_g$ 能级被降低：$\mathrm{d}_{xz}$、$\mathrm{d}_{yz}$ 能级下降所省的能量恰被
$\mathrm{d}_{xy}$ 能级的升高所抵偿。}
\end{figure}

\subsection{扬--特勒效应}

到目前为止我们假定：只需弄清局域环境具有何种对称性，推断出电子结构，再按填充
能级的电子数推出磁性质。然而有时磁性质本身竟能反过来影响局域环境的对称性！原因
在于：如八面体这样的结构自发畸变（见图 3.6）有时在能量上有利——弹性能量的增加
恰好被畸变带来的电子能量节省所平衡。这一现象称为扬--特勒效应（Jahn--Teller
effect）。例如八面体环境中的 $\mathrm{Mn^{3+}}$ 离子（组态 $3d^{4}$）就表现出这种
行为（见图 3.6）；相反，$\mathrm{Mn^{4+}}$ 离子（$3d^{3}$）不表现出，因为畸变
没有电子能量的净降低。

为在唯象层面描述这一效应，设系统的畸变可用参数 Q 量化——沿某合适简正模坐标的
畸变距离。它带来一个二次型的能量代价：
\begin{equation*}
E(Q) = \frac{1}{2}M\omega^{2}Q^{2},\tag{3.3}
\end{equation*}
M 与 $\omega$ 分别是阴离子质量和该简正模对应的角频率。此关系绘于图 3.7(a)。显然
最小畸变能为零，出现在 Q=0（无畸变）处。

畸变还会提升某些轨道的能量、降低另一些的能量。若所有轨道或全满或全空，这无关
紧要——总能量就是式 3.3。但对部分填充的轨道，这个效应可能极其重要：系统可以实现
总能量的净降低。电子能量对 Q 的依赖可能相当复杂，但可写成对 Q 的泰勒级数；只要
畸变小，只保留线性项是合理的。设某轨道的能量含 $+AQ$ 或 $-AQ$ 项（对应电子能量的
升高或降低），A 为适当的正常数。总能量 $E(Q)$ 是电子能量与弹性能量之和：
\begin{equation*}
E(Q) = \pm AQ + \frac{1}{2}M\omega^{2}Q^{2},\tag{3.4}
\end{equation*}
$AQ$ 项符号的两种选择给出两条曲线，绘于图 3.7(b)。只看其中一条，用
$\partial E/\partial Q=0$ 求该轨道的最小能量，得
\begin{equation*}
Q_0 = \frac{A}{M\omega^{2}}\tag{3.5}
\end{equation*}
及最小能量 $E_{\min}=-A^{2}/2M\omega^{2}<0$。若只有该轨道被占据，系统便可因自发
畸变而净省能量。

\begin{figure}
\centering
\includegraphics[width=0.55\textwidth]{fig3_07.pdf}
\caption{（a）按式 3.3，八面体配合物能量对畸变 Q 的函数。（b）按式 3.4 的情形。}
\end{figure}

迄今讨论的实质上是静态扬--特勒效应：自发畸变固定在八面体的某一特定轴上。但在
较高温度下，畸变可以在不同轴之间切换，产生动态扬--特勒效应。另一类动态效应是
畸变在格点之间的快速跳迁——在例如含 $\mathrm{Mn^{3+}}$ 与 $\mathrm{Mn^{4+}}$
混合离子的材料中很重要。这些现象可以通过它们对磁共振数据的影响来探测。

还有一类效应：某些材料（如 $\mathrm{DyVO_4}$）在低于某临界温度时，各配合物的
扬--特勒畸变会在整个晶体中协同发生，称为协同扬--特勒转变（cooperative Jahn--Teller
transition）。

晶体场的讨论至此告一段落。本章余下部分描述可用于研究磁矩与其环境相互作用的
各种实验技术。

\section{磁共振技术}

磁矩的局域环境由晶体场决定，但自旋--轨道耦合以及与核的超精细相互作用对电子结构
同样重要。这些效应可以用多种磁共振实验技术研究，本章余下部分将描述这些技术。

1.1.2 节已经看到：对磁矩施加磁场可诱发该磁矩以角频率 $|\gamma B|$ 进动，$\gamma$
为旋磁比。磁场中的磁矩系统可在此频率吸收能量，于是可以观察到系统对调到正确频率的
电磁波的共振吸收。这就是磁共振（magnetic resonance）；依共振磁矩的类型不同，实验
形式多样。电磁波的适当频率取决于磁场与磁矩的大小。图 3.8 以各种单位总结了电磁辐
射的频率范围，供下文随时参照。

\begin{figure}
\centering
\includegraphics[width=0.9\textwidth]{fig3_08.pdf}
\caption{电磁波谱。光子能量以温度 $T=E/k_{\mathrm{B}}$（K）与能量 $E$（eV）标出。
相应频率 $f$ 以 Hz 标出，也以光谱学常用的 $\mathrm{cm}^{-1}$ 标出。$\mathrm{cm}^{-1}$
标尺上标了若干常见分子跃迁与激发（分子转动与振动的典型范围，以及 C--H 弯曲与伸缩
模式），典型 $\pi$ 与 $\sigma$ 键的能量亦然。光子波长 $\lambda=c/f$（$c$ 为光速）。
温度标尺上特地标出 $T_{\mathrm{CMB}}$（宇宙微波背景温度）、常压下液氦（$\mathrm{He^4}$）
与液氮（$\mathrm{N_2}$）的沸点以及室温。图中其他缩写：IR=红外，UV=紫外，R=红，
G=绿，V=紫。字母 H 标记 13.6 eV——氢 1s 电子能量的大小。频率轴上还标出电磁波谱
主要区域的名称：radio（无线电）、microwave（微波）、infrared（红外，含“近”与
“远”）、optical（光学）与 UV（紫外）。}
\end{figure}

\subsection{核磁共振}

最常用的磁共振形式是核磁共振（nuclear magnetic resonance，NMR）。这一技术在医学
成像中大量使用，为避讳“核”这个吓人的字眼而称磁共振成像（MRI）。有人得知人体
平均含有超过 $10^{27}$ 个核可能会大惊失色：按质量计，我们 $\sim$99.98\% 是核！
MRI 测量的是病人体内质子的 NMR，提供安全、无创的精细断层成像。

做任何 NMR 实验都需要一个自旋非零的核。常研究的核包括 ${}^{1}$H（质子）、
${}^{2}$H（氘核）与 ${}^{13}$C。简单 NMR 实验（见图 3.9）把样品放进一个线圈内，
线圈安装在磁体极靴之间。磁体产生沿某方向（如 $z$ 方向）的磁场 $B_0$。我们已经
知道核角动量的 $z$ 分量 $m_I$ 只能取 $-I$ 到 $I$ 之间的整数间隔值。

\begin{figure}
\centering
\includegraphics[width=0.45\textwidth]{fig3_09.pdf}
\caption{NMR 实验示意图。样品置于射频（RF）线圈内，线圈产生振荡的射频场；磁体提供
高度均匀的静磁场。静场 $\mathbf{B}_0$ 与振荡场 $\mathbf{B}_1$ 互相垂直。真实实验中
样品比图中所示小得多，从而感受到 RF 线圈的均匀场。}
\end{figure}

核的能量是磁矩\fn{3}{核带正电、电子带负电，这正是下列表达式中符号不同的原因：$\boldsymbol{\mu}=+g_{\mathrm{N}}\mu_{\mathrm{N}}\mathbf{I}$（核）与 $\boldsymbol{\mu}=-g_J\mu_{\mathrm{B}}\mathbf{J}$（电子）。}$\boldsymbol{\mu}$ 在磁场 $B_0$ 中的能量：$E=-\boldsymbol{\mu}\cdot\mathbf{B}_0$，故
\begin{equation*}
E = -g_{\mathrm{N}}\mu_{\mathrm{N}}m_IB_0.\tag{3.6}
\end{equation*}
这是 $2I+1$ 个等间距能级的阶梯，间距 $g_{\mathrm{N}}\mu_{\mathrm{N}}B_0$；Na 或 Cu
中 $I=\frac{3}{2}$ 的情形示于图 3.10(a)。\fn{4}{此处先设 J=0，从而忽略电子自旋。}用射频场激发相邻能级对之间的跃迁，正是核磁共振的
基础。RF 场 $B_1$ 沿 $x$ 方向施加，对系统的微扰正比于 $B_1\hat{I}_x$。微扰矩阵元
正比于 $\langle m_I'|\hat{I}_x|m_I\rangle$，除非 $m_I'=m_I\pm1$ 否则为零。故允许
跃迁由选择定则
\begin{equation*}
\Delta m_I = \pm 1,\tag{3.7}
\end{equation*}
描述：只有相邻态之间的跃迁可以发生。最常见的装置中，RF 线圈不仅产生激发，还自身
是高品质因数（大 Q 值）调谐回路的一部分。RF 场激发核跃迁时，能量在 RF 回路与样品
之间传递，导致回路 Q 值的微小变化。

\begin{figure}
\centering
\includegraphics[width=0.45\textwidth]{fig3_10.pdf}
\caption{（a）$I=\frac{3}{2}$ 核的四个能级。（b）较简单的 $I=\frac{1}{2}$ 核只有两个
能级，间隔 $\Delta E=g_{\mathrm{N}}\mu_{\mathrm{N}}B_0$；下（上）能级对应核磁矩平行
（反平行）于磁场 $B_0$。}
\end{figure}

实验有两种做法：固定 RF 场频率而改变磁场，或固定磁场而改变 RF 频率。通常采用
前者。关键是要有高度均匀的磁体来产生恒定场 $B_0$；否则样品不同部分处于略不同的
磁场中，在磁场扫描的不同点先后共振，使测得的共振极宽，甚至可能完全被抹平。

\begin{example}{3.2}
对质子 $I=\frac{1}{2}$，$m_I$ 只能取 $\pm\frac{1}{2}$（图 3.10(b)）。系统的两态
间隔 $\Delta E=g_{\mathrm{N}}\mu_{\mathrm{N}}B$，非常小：典型实验室磁场
$B_0\sim1$~T 中质子的能隙 $\sim10^{-7}$ eV；写成 $k_{\mathrm{B}}T$ 只相当于
$\sim1$ mK 的温度。

因此室温下、该磁场处，核平均而言只有极微弱的沿外场排齐的趋势——热随机化能量远
超微弱的排齐能量。若不采用共振技术，核磁性引起的任何效应实际上都无法探测。NMR
以恰好\fn{5}{其实不必完全精确：可以差一个线宽。线宽由多种因素决定，包括相邻核的偶极场造成的核处磁场变化。}合适的角频率 $\omega$ 的 RF 信号微扰系统：
\begin{equation*}
\hbar\omega = \Delta E\tag{3.8}
\end{equation*}
从而可以激发跨越这一能隙的跃迁。为什么需要射频？下例说明。
\end{example}

\begin{example}{3.3}
磁场 1 T 时，频率为
\begin{equation*}
\nu = \frac{\omega}{2\pi} = \frac{g_{\mathrm{N}}\mu_{\mathrm{N}}B_0}{2\pi\hbar} = 42.58\ \mathrm{MHz}\tag{3.9}
\end{equation*}
处于电磁波谱的射频区。
\end{example}

这意味着可以使用振荡器与线圈（相比之下，电子自旋共振实验须用 Gunn 二极管与微波
波导，见 3.2.2 节）。于是，在恰好合适的 RF 激发频率处系统将吸收能量——这就是核磁
共振：RF 激发频率与核能级间隔之间的共振。现代 NMR 谱仪使用 12--15 T 的磁体，RF
频率在 $\sim$500--650 MHz 范围，仍在射频区。

给定磁场下 NMR 共振频率可以从 $\gamma B/2\pi$ 略微上下偏移，取决于核的化学环境。
这些化学位移（chemical shifts，典型为百万分之几）源于绕核运动的电子对核的轻微
屏蔽。给定化学环境中核被屏蔽的程度是已知的，因此可以对特定分子“指纹化”。NMR
谱线还会因与相邻核的磁耦合而分裂，称为自旋--自旋耦合（spin--spin coupling），经由
电子传递的间接接触超精细相互作用发生。这同样给出核环境的信息。

为更细致地理解我们诱导的跃迁，考虑双能级自旋系统（$I=\frac{1}{2}$）如何从 RF
线圈吸收能量。记下能级为 $-$、上能级为 $+$。诱导两能级间跃迁的概率与跃迁方向无关，
因为振荡 RF 场对应的微扰哈密顿量 $\hat{H}\propto\hat{I}_x$ 是厄米的：
\begin{equation*}
|\langle\psi_-|\hat{\mathcal{H}}|\psi_+\rangle|^{2} = |\langle\psi_+|\hat{\mathcal{H}}|\psi_-\rangle|^{2}.\tag{3.10}
\end{equation*}
于是所谓受激跃迁（+ 与 $-$ 之间）的概率密度与方向无关，以速率 W 发生，W 正比于
所用 RF 功率。设 $t$ 时刻下能级有 $N_-(t)$ 个自旋，则每单位时间 $WN_-(t)$ 个被
激发到上能级；上能级有 $N_+(t)$ 个时，每单位时间 $WN_+(t)$ 个被激发到下能级。故
\begin{equation*}
\frac{\mathrm{d}N_+(t)}{\mathrm{d}t} = WN_-(t) - WN_+(t)\tag{3.11}
\end{equation*}
\begin{equation*}
\frac{\mathrm{d}N_-(t)}{\mathrm{d}t} = WN_+(t) - WN_-(t)\tag{3.12}
\end{equation*}
相减得
\begin{equation*}
\frac{\mathrm{d}}{\mathrm{d}t}(N_+(t)-N_-(t)) = -2W(N_+(t)-N_-(t)),\tag{3.13}
\end{equation*}
解为
\begin{equation*}
N_+(t)-N_-(t) = (N_+(0)-N_-(0))\,\mathrm{e}^{-2Wt}\tag{3.14}
\end{equation*}
初始的布居差在受激电磁跃迁驱动下指数趋于零。定义
\begin{equation*}
n(t) = N_+(t)-N_-(t),\tag{3.15}
\end{equation*}
式 3.14 成为
\begin{equation*}
n(t) = n(0)\,\mathrm{e}^{-2Wt}\tag{3.16}
\end{equation*}

自旋系统吸收或发射电磁能量的速率容易计算：每次跃迁交换能量 $\hbar\omega$。系统
在 $t$ 时刻的能量为
\begin{equation*}
E(t) = N_-E_- + N_+(E_-+\hbar\omega),\tag{3.17}
\end{equation*}
$E_-$ 是下能级能量。此式可整理成常数项加 $\frac{1}{2}\hbar\omega\,n(t)$。于是能量
吸收速率为
\begin{equation*}
\frac{\mathrm{d}E}{\mathrm{d}t} = -W\hbar\omega\,n(t)\tag{3.18}
\end{equation*}
它与 $n(t)$ 成正比，随上、下能级布居逐渐相等而以时间常数 $1/2W$ 趋于零。这也表明
吸收能量需要 $n(t)\neq0$，即布居差。布居差由系统与热模式的相互作用产生，使系统
的极化回到玻尔兹曼概率：
\begin{equation*}
\left(\frac{N_+}{N_-}\right)_0 = \mathrm{e}^{-\hbar\omega/k_{\mathrm{B}}T}\tag{3.19}
\end{equation*}
下标 0 表示系统的平衡值。吸收能量后极化可能改变；若核自旋系统与样品的热运动及
激发相互作用，自旋系统的极化可以回到平衡值。若自旋系统的极化被快速的电磁跃迁
降为零，关闭电磁跃迁后自旋系统需要一段时间“恢复”；这段时间称为 $T_1$，度量
自旋系统与其环境相互作用的时间常数。$T_1$ 是自旋--晶格弛豫时间（spin--lattice
relaxation time，也称纵向弛豫时间）。预期极化按
\begin{equation*}
n(t) = n_0\left(1-\mathrm{e}^{-t/T_1}\right).\tag{3.20}
\end{equation*}
恢复。

\begin{figure}
\centering
\includegraphics[width=0.5\textwidth]{fig3_11.pdf}
\caption{式 3.25 给出的电磁能量吸收速率随跃迁速率 W 的变化。W 小时正比于 W；W 大时
饱和于由 $1/T_1$ 决定的值。}
\end{figure}

现在把两个过程——受激跃迁与弛豫——放在一起。合并方程得
\begin{equation*}
\frac{\mathrm{d}}{\mathrm{d}t}n(t) = -2Wn(t) + \frac{n_0-n(t)}{T_1},\tag{3.21}
\end{equation*}
在
\begin{equation*}
\frac{\mathrm{d}}{\mathrm{d}t}n(t) = 0.\tag{3.22}
\end{equation*}
时有稳态解
\begin{equation*}
n(t) = \frac{n_0}{1+2WT_1},\tag{3.23}
\end{equation*}
电磁能量的稳态吸收速率为
\begin{align*}
\frac{\mathrm{d}E}{\mathrm{d}t} &= n(t)\hbar\omega W\tag{3.24}\\
&= n_0\hbar\omega\frac{W}{1+2WT_1}.\tag{3.25}
\end{align*}
此关系绘于图 3.11。低幅 RF 微扰场下，能量吸收速率正比于 W；大幅 RF 微扰场下，
它稳定在与 $T_1$ 成正比、与 W 的精确大小无关的水平。这称为饱和（saturation）。

电磁能量吸收速率正比于上、下能级的布居差，因此可用它监测布居差及其时间依赖——
当然这种观测本身必然在一定程度上扰动布居差。式 3.25 与图 3.11 表明：提高 RF 功率
只能到某个程度为止。式 3.25 还表明 $\mathrm{d}E/\mathrm{d}t$ 正比于 $\hbar\omega$：
能量更高（频率更高）的光子给出更大信号；但这当然意味着需要更大的磁场使 $g\mu_{\mathrm{N}}B$
匹配 $\hbar\omega$。这对 ESR（见下节 3.2.2）反而是优势：激发跨越更大的电子能隙
需要很强健的微波光子，从而带来更高的灵敏度。由于电子磁矩大于核磁矩，ESR 即便用
比 NMR 更小的磁场，也仍能使用频率更高的光子。
\bio{Felix Bloch\\（1905--1983）}

回到 NMR，更细致地考虑弛豫过程。磁场中自旋存在很弱的极化（弱是因为能隙
$\Delta E\ll k_{\mathrm{B}}T$）。设接通静磁场 $B_0$ 后（回顾 $B_0$ 沿 $z$ 轴）
极化取平衡值 $M_0$。射频激发的效果是逐渐破坏这一磁化强度；而如前所述，切断磁场
后，磁化强度通过与周围的弱相互作用以时间常数 $T_1$ 弛豫回平衡值。因此预期
\begin{equation*}
\frac{\mathrm{d}M_z}{\mathrm{d}t} = \frac{M_0-M_z}{T_1}.\tag{3.26}
\end{equation*}
$T_1$ 弛豫必须涉及与晶格的相互作用，因为能量必须与晶格交换：改变 $M_z$ 有能量
后果——自旋沿 $z$ 方向弛豫会改变它相对外加场的取向，从而改变其能量。

$M_x$ 与 $M_y$ 分量应为零；若不为零，将以时间常数 $T_2$ 弛豫回零：
\begin{equation*}
\frac{\mathrm{d}M_x}{\mathrm{d}t} = -\frac{M_x}{T_2},\qquad
\frac{\mathrm{d}M_y}{\mathrm{d}t} = -\frac{M_y}{T_2}.\tag{3.27}
\end{equation*}
$T_2$ 是自旋--自旋弛豫时间，是失相（dephasing）的特征时间，对应自旋系统不同部分
之间的相互作用；也可能由 $B_0$ 的不均匀性造成。它导致被观测自旋与邻近自旋相互
作用产生的进动频率差异。改变 $M_x$ 或 $M_y$ 没有能量后果，因为外加场沿 $M_z$。

外加磁场还引起自旋进动，于是 $M_x$、$M_y$、$M_z$ 的方程为
\begin{equation*}
\frac{\mathrm{d}M_x}{\mathrm{d}t} = \gamma(\mathbf{M}\times\mathbf{B})_x - \frac{M_x}{T_2},\tag{3.28}
\end{equation*}
\begin{equation*}
\frac{\mathrm{d}M_y}{\mathrm{d}t} = \gamma(\mathbf{M}\times\mathbf{B})_y - \frac{M_y}{T_2},\tag{3.29}
\end{equation*}
\begin{equation*}
\frac{\mathrm{d}M_z}{\mathrm{d}t} = \gamma(\mathbf{M}\times\mathbf{B})_z + \frac{M_0-M_z}{T_1},\tag{3.30}
\end{equation*}
称为布洛赫方程（Bloch equations）。它们于 1946 年导出，可在许多有意义的情形下
求解，常用旋转参照系法——把坐标换成在 $xy$ 平面内以共振频率旋转的坐标系。

自旋--自旋弛豫时间 $T_2$ 可用如下技术测量：在与 $B_0$ 热平衡、$B_1$ 关断时，
存在沿 $z$ 的弱磁化强度。加一个时长 $t_{\mathrm{p}}$ 的短 RF 脉冲（频率调到接近
自旋共振频率），使自旋转过角度 $\gamma B_1t_{\mathrm{p}}$（$B_1$ 为 RF 振幅）。
角度可通过改变 $t_{\mathrm{p}}$ 调节；取 $t_{\mathrm{p}}=\pi/2\gamma B_1$ 就得到
所谓的 $90^{\circ}$ 脉冲。脉冲幅度与时长选得使磁化强度恰好进动一小段时间、转入
$xy$ 平面。关断 $B_1$ 后，磁化强度在静场 $\mathbf{B}_0$ 中进动，在 $xy$ 平面内以
频率 $\gamma B_0$ 稳定旋转，在线圈中产生感应电压。若没有 $T_2$ 弛豫过程（自旋--
自旋弛豫），振荡将永远持续；正是 $T_2$ 使振荡衰减，从而给出 $T_2$ 的测量。线圈
电压的相应衰减称为自由感应衰减（free induction decay，见图 3.12）。

\begin{figure}
\centering
\includegraphics[width=0.55\textwidth]{fig3_12.pdf}
\caption{模拟的自由感应衰减信号。振荡因自旋--自旋弛豫（$T_2$ 过程）随时间衰减。
数据呈现干涉图样，来自三个不等价核——它们因在晶体中所处格位不同而感受略微不同的
磁场。}
\end{figure}

这一方法的问题在于它同时测到了磁体不均匀性导致的弛豫——样品不同部位进动频率
略异。需要一种方法把真正的自旋--自旋弛豫与磁体不均匀性弛豫分开。办法之一是自旋
回波（spin echo）技术，其原理示于图 3.13：先用 $90^{\circ}$ 脉冲（如前把自旋转入
$xy$ 平面），随后加 $180^{\circ}$ 脉冲（把自旋转过 $180^{\circ}$）。可以把这一
效应比作一场马拉松：所有选手稳步跑，但有的快有的慢。发令枪（$90^{\circ}$ 脉冲）
后选手们成群出发，但不久体力好的渐渐领先，其余的人渐渐落后。然后突然在时刻 $\tau$，
规则改变（$180^{\circ}$ 脉冲）：选手们被告知掉头跑回起点。此时慢者反而占优——
他们离起点更近。由对称性易见：无论各自速度如何，所有选手将在时刻 $2\tau$ 同时
回到起点。同理，所有自旋在 $2\tau$ 后重新对齐，不论各自的失相。这就是自旋回波。

\begin{figure}
\centering
\includegraphics[width=0.46\textwidth]{fig3_13.pdf}
\caption{自旋回波效应。（a）平衡磁化强度初始沿 $z$ 方向、平行于 $B_0$。（b）称该时刻为 $t=0$：沿 $x$ 轴的 $90^{\circ}$ 脉冲把自旋转入 $xy$ 平面。（c）由于 $z$ 轴方向的恒定场 $B_0$，自旋在 $xy$ 平面内进动，但因磁场不均匀而速率略异。（d）它们逐渐相互失相。（e）$t=\tau$ 时沿 $x$ 轴的 $180^{\circ}$ 脉冲使自旋绕 $x$ 轴转过 $180^{\circ}$；随后在 $B_0$ 中进动时次序已颠倒。（f）$2\tau$ 时它们重新聚齐，产生自旋回波信号——前提是期间没有发生其他弛豫过程。图中画的是两脉冲间隔较短的情形；若间隔长、自旋在 $180^{\circ}$ 脉冲前已转过若干圈，效应同样成立。}
\end{figure}

于是实验中可用自旋回波消除磁场不均匀导致自旋以不同速率进动的问题，化学位移造成
的场展宽也被消除。但自旋--自旋相互作用——来自相邻核的含时涨落随机磁场——的效应
无法消除：这一弛豫机制无法重聚焦，回波幅度将随 $\tau$ 衰减。回波信号的 NMR 强度
应遵循 $I(2\tau)=I(0)e^{-2\tau/T_2}$，由此可测得真正的（自旋--自旋耦合导致的）
$T_2$。

$T_1$ 可用类似方法测量，只是先加 $180^{\circ}$ 脉冲：它使磁化强度从 $+\hat{\mathbf{z}}$
转到 $-\hat{\mathbf{z}}$，随后以时间常数 $T_1$ 弛豫回 $+\hat{\mathbf{z}}$，但不
进动。于是 $180^{\circ}$ 脉冲后 $\tau$ 时刻
\begin{equation*}
\mathbf{M}(\tau) = \hat{\mathbf{z}}\,M_z(\tau) = \hat{\mathbf{z}}\,M_0(1-2\,\mathrm{e}^{-\tau/T_1}).\tag{3.31}
\end{equation*}
$180^{\circ}$ 脉冲后 $\tau$ 时刻再用 $90^{\circ}$ 脉冲把磁化强度转入 $xy$ 平面，
开始初始幅度为 $M_z(\tau)$ 的自由感应衰减。于是测量自由感应衰减初始幅度对
$180^{\circ}$--$90^{\circ}$ 脉冲间隔时间的依赖，即可得 $T_1$。其他多种脉冲序列
也可用于提取更多类型的信息。

\subsection{电子自旋共振}

可以用电子做与 NMR 类似的实验，称为电子自旋共振（electron spin resonance，ESR），
有时也称电子顺磁共振（electron paramagnetic resonance，EPR）。电子磁矩远大于核
磁矩，进动频率也高得多：典型实验室磁场下，电磁辐射已处于微波波段。实验装置示于
图 3.14：样品置于谐振腔内，微波经波导进入。通常更方便的做法是固定微波频率、扫描
磁场。谐振腔品质因数极高，用以提高微弱 ESR 信号的探测灵敏度。记录微波吸收随外加
磁场的变化。常在磁体提供的静场之外加一组调制线圈，在样品处产生小的振荡磁场，
再配合相敏检测技术进一步提高测量灵敏度。\fn{6}{锁相放大器产生驱动振荡磁场的信号；微波信号馈入锁相放大器，后者滤除任何不在振荡磁场频率处的成分，从而大幅降低噪声。}

\begin{figure}
\centering
\includegraphics[width=0.45\textwidth]{fig3_14.pdf}
\caption{ESR 实验示意图。微波经波导进入谐振腔，共振诱导的微波吸收通过监测腔的
Q 因子来测量。样品必须放在磁体中心——那里场最均匀。}
\end{figure}

ESR 跃迁的选择定则是
\begin{equation*}
\Delta m_J = \pm 1,\tag{3.32}
\end{equation*}
即熟悉的偶极选择定则——我们诱导的是磁偶极跃迁。图 3.15 展示八面体环境中
$J=1$ 离子（如 $\mathrm{Ni^{2+}}$）的这些跃迁。实验在固定频率 $h\nu$ 下扫描磁场：
每当两个相邻能级间隔为 $h\nu$ 就发生一次 ESR 跃迁。图 3.15 表明谱线的位置与数目
可以给出晶体场劈裂的信息。

\begin{figure}
\centering
\includegraphics[width=0.45\textwidth]{fig3_15.pdf}
\caption{$J=1$ 离子（如 $\mathrm{Ni^{2+}}$）的 ESR：（a）无晶体场劈裂（单条 ESR
线）；（b）有晶体场劈裂 $\Delta$（两条 ESR 线）。}
\end{figure}

图 3.16 给出自由 $\mathrm{Mn^{2+}}$ 离子（自旋 $\frac{5}{2}$）的能量随磁场的变化。
磁场 B 下不同 $m_J$ 能级间有许多可能跃迁，大小都是 $g\mu_{\mathrm{B}}B$。因此
ESR 实验或许预期在 $h\nu=g\mu_{\mathrm{B}}B$ 给出的频率 $\nu$ 处观察到单线。但对
含 $\mathrm{Mn^{2+}}$ 离子的固体做真实实验时，即便在立方环境中，晶体场也会使图像
显著复杂化。

不过，核与电子之间的超精细耦合在哈密顿量中给出 $A\mathbf{I}\cdot\mathbf{J}$ 项，
每个 $m_J$ 能级按 $m_I$ 分裂成 $2I+1$ 个超精细能级。ESR 跃迁的选择定则因此是式
3.32 加上
\begin{equation*}
\Delta m_I = 0.\tag{3.33}
\end{equation*}

\begin{figure}
\centering
\includegraphics[width=0.5\textwidth]{fig3_16.pdf}
\caption{自旋 $\frac{5}{2}$ 磁矩（如自由 $\mathrm{Mn^{2+}}$ 离子）的能量随磁场的
变化。按选择定则 $\Delta m_J=\pm1$ 的可能跃迁如图。此图忽略了晶体场。}
\end{figure}

后一定则成立是因为微波只诱导电子能级间的偶极跃迁、而非核能级间的跃迁。\fn{7}{这是因为核能级间距与微波光子能量相比极小。}由于 ESR 实验的典型场下 A 远小于
$g\mu_{\mathrm{B}}B$，只需计算 J 沿 B 的分量，磁场中能级即为
\begin{equation*}
E = g\mu_{\mathrm{B}}Bm_J + Am_Im_J\tag{3.34}
\end{equation*}
（外加磁场中核自旋的能量已忽略——它比超精细能量还小得多）。用选择定则可得共振
应发生在
\begin{align*}
h\nu &= [g\mu_{\mathrm{B}}B(m_J+1) + Am_I(m_J+1)] - [g\mu_{\mathrm{B}}Bm_J + Am_Im_J]\\
&= g\mu_{\mathrm{B}}B + Am_I\tag{3.35}
\end{align*}
处，每条 ESR 线分裂成 $2I+1$ 条超精细线。

\begin{example}{3.4}
自由 $\mathrm{Mn^{2+}}$ 离子 $I=\frac{5}{2}$，故 ESR 谱中有六条线。这一情形示于
图 3.17（晶体中的 $\mathrm{Mn^{2+}}$ 离子因晶体场而更复杂：劈裂非平凡地依赖磁场
与晶轴的夹角）。
\end{example}

\begin{figure}
\centering
\includegraphics[width=0.55\textwidth]{fig3_17.pdf}
\caption{固定磁场下自由 $\mathrm{Mn^{2+}}$ 离子的超精细劈裂。$J=\frac{5}{2}$、
$I=\frac{5}{2}$，六个 $m_J$ 能级各自分裂成六个超精细能级。$\Delta m_J=\pm1$、
$\Delta m_I=0$ 的可能跃迁如图，在 ESR 谱中给出六条不同的线。此图忽略晶体场效应——
它使劈裂强烈依赖于外加磁场相对晶轴的方向。}
\end{figure}

ESR 实验在有用的范围内探测电子自旋，因为晶体场劈裂常在 GHz 频段。顺磁盐的 ESR
研究已经提供了大量关于晶体场与配体场的细致信息。效应往往是各向异性的——磁场
相对晶轴旋转时谱线会移动。

\begin{example}{3.5}
举一个如何获得这类信息的例子：考虑轴对称环境中的 $\mathrm{Ce^{3+}}$（$4f^{1}$）
（见图 3.18）。该离子 L=3、$S=\frac{1}{2}$，自旋--轨道相互作用把它分裂成八重简并
的 $J=\frac{7}{2}$ 能级和较低的六重简并的 $J=\frac{5}{2}$ 能级。晶体场是轴对称的，
按对称性只含 $J_z$ 的偶次幂，于是六重 $J=\frac{5}{2}$ 多重态（manifold，即一组谱线的合称）分裂
成三个二重态，八重 $J=\frac{7}{2}$ 多重态分裂成四个。这些劈裂在稀土体系中远小于
自旋--轨道相互作用。二重态可被磁场进一步分裂——对典型实验室磁场而言，这几乎总是
最小的能量项，通常约 $10^{-4}$ eV。低温下只有基态被占据，观察到的正是基态跃迁。
\end{example}

许多体系的基态像 $\mathrm{Ce^{3+}}$ 一样是二重态。这是克拉默斯定理（Kramers
theorem）的结论：含奇数个电子的系统在无磁场时至少保留二重简并。涉及的态对称为
克拉默斯二重态（Kramers doublets），它们互为时间共轭\fn{8}{即互为复共轭，因而互为时间反演的版本。}\bio{Hendrik A. Kramers\\（1894--1952）}——可以被磁场分裂而不能被电场分裂。$\mathrm{Ce^{3+}}$（$4f^{1}$）的 4f 壳层有奇数个电子，因此是克拉默斯离子。

$\mathrm{Ce^{3+}}$ 的基态是克拉默斯二重态；孤立地看它像自旋 $\frac{1}{2}$ 态，
尽管量子数全然不同。事实上可以赋予它有效自旋 $\tilde{S}=\frac{1}{2}$：这是虚构的
角动量，选它使所考虑的那组能级的简并度等于 $(2\tilde{S}+1)$。这样做的动机是找到
一个有效自旋哈密顿量，用某种简单得多的语言——自旋 $\frac{1}{2}$ 的自由原子或
离子——近似描述实验情形。该哈密顿量可写为
\begin{equation*}
\tilde{H} = g\mu_{\mathrm{B}}\mathbf{B}\cdot\tilde{\mathbf{S}}\tag{3.36}
\end{equation*}
这里的 $g$ 是有效 $g$ 因子（有时称光谱分裂因子）：对真实的孤立自旋
$\frac{1}{2}$ 电子它等于 2，此处还包含了轨道贡献。常常还需要用有效 $g$ 张量
$\mathbf{g}$ 表达自旋与场的相互作用依赖于磁场取向这一事实，式 3.36 须换成
\begin{equation*}
\tilde{H} = \mu_{\mathrm{B}}(\mathbf{B}\cdot\mathbf{g}\cdot\tilde{\mathbf{S}}) = \sum_{ij}\mathbf{g}_{ij}B_i\tilde{S}_j.\tag{3.37}
\end{equation*}
晶体场的效果可再往哈密顿量里加一项——表示自旋因晶体场而偏好沿特定晶向的能量。
这称为单离子各向异性（single ion anisotropy），哈密顿量中的项 $\tilde{H}_{\mathrm{SI}}$
在单轴晶体（存在特殊轴，能量只依赖自旋与该轴的夹角）为
\begin{equation*}
\tilde{H}_{\mathrm{SI}} = -DS_z^{2},\tag{3.38}
\end{equation*}
在立方晶体为
\begin{equation*}
\tilde{H}_{\mathrm{SI}} = -D(S_x^{4}+S_y^{4}+S_z^{4})\tag{3.39}
\end{equation*}
$D$ 是各向异性常数。

\begin{figure}
\centering
\includegraphics[width=0.55\textwidth]{fig3_18.pdf}
\caption{$\mathrm{Ce^{3+}}$ 离子的能级。主导的能量劈裂是自旋--轨道相互作用（此
能量约对应 3000 K，故上多重态的热布居可忽略）。两个多重态被假设的轴对称晶体场
分裂成若干二重态组；图中晶体场劈裂被夸大以便观察。较低处的劈裂 $\sim5$ meV，约
对应 50 K。最后，外加磁场分裂所有二重态（图中只画出了下面那个的分裂）。}
\end{figure}

ESR 不仅用于研究含顺磁离子盐类的物理，还广泛用于化学。共振对原子在分子中的位置
可以极为敏感，该技术被大量用于化学反应研究，尤其是自由基（free radicals）——
带未成对电子的原子或分子碎片，往往高活性，具有重要的环境与生物学意义。ESR 不如
NMR 用得广，因为它只适用于有未成对自旋的材料。

\begin{example}{3.6}
ESR 大显身手的一个领域是研究常含过渡金属离子的生物分子。一个重要例子是血红素
（haem）——铁卟啉基团，存在于血红蛋白与肌红蛋白中。这些分子参与氧及其他小分子
（如一氧化碳）的转运。血红素中的 Fe(II) 离子在与小分子结合后从 S=2 变为 S=0
（这是自旋转变的一个例子，见 3.1.1 节）。
\end{example}

ESR 中的自旋--晶格弛豫时间常常比 NMR 短得多，因为电子磁矩与晶格振动的耦合比核
强得多。晶格振动调制晶体场，经自旋--轨道相互作用耦合到电子磁矩，常使 $T_1$ 短到
ESR 谱线宽得无法观察（线宽正比于 $T_1^{-1}$）。有利的情形是 $\mathrm{Mn^{2+}}$——
它是 S 态离子（$3d^{5}$，${}^{6}S_{5/2}$），没有可供自旋耦合的轨道矩，其 ESR 在
室温下即可轻易观察。对其他过渡金属离子，$T_1$ 的大小取决于轨道淬灭的程度。\fn{9}{轨道矩被淬灭得越彻底，电子磁矩与晶格振动的耦合越小，$T_1$ 越长，ESR 线越窄。}冷却样品可使共振变窄（$T_1$ 增大），因为晶格振动减少；这还提高了能级间相对热
布居，增强共振强度。

ESR 的一个变种称为 ENDOR（电子核双共振，Electron Nuclear Double Resonance 的
缩写），可以以极高精度测量超精细相互作用常数与核塞曼相互作用。顺磁盐中核能级间
的 NMR 跃迁往往太弱而无法直接测量。ENDOR 同时使用 RF 信号与微波信号，借力 ESR 的
高灵敏度：微波功率调到 $\Delta m_I=0$ 的跃迁，此时 ESR 信号依赖于相应能级的相对
布居；再用调到 NMR 跃迁的 RF 功率改变布居，NMR 跃迁便可经由它对 ESR 信号的影响
被探测。

\subsection{穆斯堡尔谱学}

另一类谱学基于 1957 年发现的穆斯堡尔效应（M\"ossbauer effect）。实验原理示于
图 3.19。含 ${}^{57}$Co 核的源提供现成的激发态 ${}^{57}$Fe 核；它们经级联 $\gamma$
跃迁衰变到基态，其中包含一条 14.4 keV 的 $\gamma$ 射线（对应频率
$\nu=3.5\times10^{18}$ Hz）。若这条 $\gamma$ 射线被共振吸收，就能激发所研究样品中
的跃迁——其能量须匹配样品中的能隙。让源以速度 $v$ 运动，可因多普勒效应非常轻微地
调节 $\gamma$ 射线频率。光子频率极高，多普勒频移相当可观：$v=1\ \mathrm{mm\,s^{-1}}$
导致 $\nu v/c\approx12$ MHz 的频移。于是可以探测源或吸收体核中可能由磁相互作用或
其他相互作用引起的基态劈裂。

\begin{figure}
\centering
\includegraphics[width=0.55\textwidth]{fig3_19.pdf}
\caption{穆斯堡尔技术原理示意。${}^{57}$Co 缓慢衰变到 ${}^{57}$Fe 核的激发态；
随后 91\% 的衰变迅速到达 $I=\frac{3}{2}$ 态，再衰变到 $I=\frac{1}{2}$ 基态，放出
14.4 keV 光子。实验中让 ${}^{57}$Co 源相对样品以速度 $v$ 运动（样品须含一些
${}^{57}$Fe 核）。探测器测量 $\gamma$ 射线穿过样品的透射，由此推知吸收。}
\end{figure}

倘若激发态 ${}^{57}$Fe 核发射的 $\gamma$ 光子频率不明确，这项技术将毫无用处。
这正是 ${}^{57}$Fe 核相对缓慢的衰变至关重要的原因：0.2 $\mu$s 的半衰期只对应
约 2 MHz 的频率不确定度。第二个关键特征是 ${}^{57}$Fe 原子处于固态。为守恒动量，
自由 Fe 原子会获得 $h\nu/m_{\mathrm{Fe}}c\sim80\ \mathrm{m\,s^{-1}}$ 量级的反冲
速度，实验就毁了（我们已见过 1 mm\,s$^{-1}$ 的相对速度即可产生可测效应）。然而
被刚性固定在固体中的 ${}^{57}$Fe 原子把动量传递给整个晶体：若反冲能低于某个量，
声子发射的概率趋于零，$\gamma$ 射线可以不损失任何反冲能量地发射。这种无反冲
发射与共振吸收正是穆斯堡尔效应的精髓。

上述条件意味着：低能 $\gamma$ 射线 + 强结合于晶格的核 + 低温，效应最佳。
${}^{57}$Fe 并非唯一合适的同位素（可用者还有 ${}^{119}$Sn、${}^{127}$I、
${}^{151}$Eu、${}^{197}$Au），但最常用。穆斯堡尔测量的原始结果是一张 $\gamma$
计数（或相对吸收）对源与吸收体相对速度的图。

这个效应为什么能告诉我们样品的信息？首先，共振吸收未必出现在“期望”的位置
（源静止时），而可能略有偏移（即源以某特定速度运动时）。这一同质异能移（isomer
shift）源于核体积上核电荷分布与电子电荷分布之间库仑相互作用的微小变化——与
${}^{57}$Fe 核在 $I=\frac{3}{2}$ 态体积略增有关。

此外，观察到的共振吸收线可能不止一条，而是若干条。原因可以是四极劈裂（quadrupole
splitting）或磁劈裂（magnetic splitting）。前者源于激发态 ${}^{57}$Fe 核的电
四极矩（${}^{57}$Fe 基态 $I=\frac{1}{2}$ 无电四极矩，但所关心的激发态
$I=\frac{3}{2}$，$I>\frac{1}{2}$ 的核可有非零四极矩）。若核处于电场梯度中——某些
晶体环境即是如此——核四极矩与电场梯度的相互作用把激发的 $I=\frac{3}{2}$ 态分裂成
二重态，穆斯堡尔谱中产生两条线。后者由核与局域磁场的相互作用造成：它把 $I=\frac{1}{2}$
基态分裂成二重态、把激发的 $I=\frac{3}{2}$ 态分裂成四重态，穆斯堡尔谱中可有六条线
（选择定则 $\Delta m_I=0,\pm1$）。这一效应可用来探测磁交换相互作用与局域磁场。
所有这些移动与劈裂示意于图 3.20。注意图中（也绝无法）按比例画：观察到的劈裂典型
在 $10^{7}$--$10^{8}$ Hz，比 ${}^{57}$Fe 核基态（$I=\frac{1}{2}$）与激发态
（$I=\frac{3}{2}$）之间的能隙小 11 或 12 个量级。换句话说：这是在用 14.4 keV 的
光子测量小于 1 $\mu$eV 的能隙！若非穆斯堡尔效应——这种无反冲的共振吸收与发射——
此类实验完全不可能。

\begin{figure}
\centering
\includegraphics[width=0.55\textwidth]{fig3_20.pdf}
\caption{化学位移、四极劈裂与磁劈裂对 ${}^{57}$Fe 核能级的影响。箭头表示穆斯堡尔
吸收跃迁。跃迁大小的差异被极大夸张：实际上它们彼此之差小于 $10^{11}$ 分之一。}
\end{figure}

\subsection{$\mu$ 子自旋旋转}

$\mu$ 子是自旋 $\frac{1}{2}$ 粒子（电荷 $\pm e$，质量 $250m_{\mathrm{e}}$），寿命
$2.2\ \mu$s。在自然界，$\mu$ 子是到达海平面的宇宙线的主要成分。$\mu$ 子可用于
研究样品的磁性质，但宇宙线源强不足，须用同步加速器与回旋加速器提供的更强 $\mu$
子束。

下面较详细地描述称为 $\mu$ 子自旋旋转（muon-spin rotation，常缩写为 $\mu$SR）
的技术。必须认识到：与本书后面讨论的中子与 X 射线技术截然不同，这里不涉及散射；
$\mu$ 子被注入感兴趣的样品并在其短暂余生中留在那里、永不再出来。从样品中释放
出来并给出 $\mu$ 子信息的是它们衰变成的正电子。

$\mu$ 子质量介于电子与质子之间，其磁矩与旋磁比亦然。$\mu$ 子有正负两种电荷态，
磁性实验中特别有用的是正 $\mu$ 子 $\mu^{+}$。作为小的带正电粒子，它被电子密度大
的区域吸引，停在无机材料的间隙位，或直接键连到有机分子上。相反，负 $\mu$ 子
$\mu^{-}$ 注入后停在原子核附近，对磁性质一般远不敏感。

用高能质子束轰击合适的靶产生 $\pi$ 介子，$\pi$ 介子（$\pi^{+}$）极快地（26 ns）
衰变为 $\mu$ 子：
\begin{equation*}
\pi^{+} \rightarrow \mu^{+} + \nu_{\mu},
\end{equation*}
若只选取在靶中停下来的 $\pi$ 介子所产生的 $\mu$ 子，则 $\mu$ 子束完全自旋极化。
原因在于：同时产生的中微子（$\nu_{\mu}$）自旋与其动量反平行，而 $\pi$ 介子无自旋，
故 $\mu$ 子自旋必须与其动量反平行——出射即自旋极化。这些 $\mu$ 子随后被注入样品，
但其能量不小，至少 4 MeV。注入后它们通过原子电离与电子散射，在 0.1--1 ns 内很快
损失能量到几个 keV；接着经过一连串电子俘获与丢失反应，在约一皮秒内降到几百 eV。
若最终形成 $\mu$ 子素（muonium，见图 3.21）——$\mu^{+}$ 与 $e^{-}$ 组成的类氢态——
电子俘获最终胜出，最后几个 eV 通过 $\mu$ 子素原子与宿主原子的非弹性碰撞损失。
这一切都极快，$\mu$ 子（或 $\mu$ 子素）迅速热化。而且这些效应全源于库仑作用，
不与 $\mu$ 子自旋相互作用，故 $\mu$ 子在物质中热化时几乎不退极化——这对 $\mu$SR
实验至关重要。有人可能担心 $\mu$ 子只测量到被高能入射 $\mu$ 子辐射损伤的区域。
这实际上不构成问题：空位产生有阈值能量，只有 $\mu$ 子路径的初始部分受损严重；
过了损伤点，$\mu$ 子仍有足够能量继续在样品中前行约 1 $\mu$m，远离任何诱发产生的
空位与间隙原子。

\begin{figure}
\centering
\includegraphics[width=0.4\textwidth]{fig3_21.pdf}
\caption{$\mu$ 子素：由 $\mu^{+}$ 与 $e^{-}$ 组成的类氢态。}
\end{figure}

$\mu$SR 实验中，$\mu$ 子停在感兴趣的样品里，经时间 $t$ 后以正比于
$\mathrm{e}^{-t/\tau_{\mu}}$ 的概率衰变，$\tau_{\mu}=2.2\ \mu$s 是 $\mu$ 子寿命。
$\mu$ 子衰变是三体过程
\begin{equation*}
\mu^{+} \rightarrow e^{+} + \nu_e + \bar{\nu}_{\mu}.
\end{equation*}
衰变经由弱相互作用，因而具有宇称不守恒这一不寻常特征。这一现象（也是 $\mu$ 子
中微子负螺旋度的根源）使发射出的正电子（$e^{+}$）倾向沿 $\mu$ 子衰变时自旋所指
的方向出射。

能量最高的正电子的角分布示于图 3.22。实际上各种能量的正电子都会发射，净效果
没有这么显著。实验中可施加垂直于初始 $\mu$ 子自旋方向的磁场，使 $\mu$ 子自旋
进动。对许多 $\mu$ 子重复实验，只要愿意采集足够长时间的数据，就能以任意精度跟踪
进动 $\mu$ 子集合的极化。实验示意于图 3.23(a)。考虑一个极化与动量反平行的 $\mu$
子被注入样品（反平行是因为它的产生方式，见上——故 $\mu$ 子进入样品时自旋指向它
来的方向）。若它“倒霉”地立即衰变，就没时间进动，正电子将优先进入后向探测器；
若多活一会儿，就有时间进动——活过半圈，正电子将优先进入前向探测器。于是，进动
$\mu$ 子集合发出的正电子束可以比作灯塔的光束。

\begin{figure}
\centering
\includegraphics[width=0.45\textwidth]{fig3_22.pdf}
\caption{发射正电子相对初始 $\mu$ 子自旋方向的角分布（能量最高的正电子的期望
分布）。}
\end{figure}

\begin{figure}
\centering
\includegraphics[width=0.55\textwidth]{fig3_23.pdf}
\caption{（a）$\mu$SR 实验示意。自旋极化的 $\mu$ 子束停在样品 S 中；衰变后，正电子
由前向探测器 F 或后向探测器 B 探测。若如图对样品施加横向磁场 H，$\mu$ 子将进动。
（b）前向（虚线）与后向（实线）探测器探测到的正电子数；点线是两信号的平均。
（c）不对称度函数。}
\end{figure}

前向与后向探测器探测到的正电子数的时间演化由函数 $N_{\mathrm{F}}(t)$ 与
$N_{\mathrm{B}}(t)$ 描述（图 3.23(b)）。由于 $\mu$ 子衰变是放射性过程，两项之和
呈指数衰减。$\mu$ 子极化的时间演化可由两者的归一化差——不对称度函数（asymmetry
function）$A(t)$——得到：
\begin{equation*}
A(t) = \frac{N_{\mathrm{B}}(t)-N_{\mathrm{F}}(t)}{N_{\mathrm{B}}(t)+N_{\mathrm{F}}(t)},\tag{3.40}
\end{equation*}
示于图 3.23(c)。$\mu$ 子产生时具有 100\% 自旋极化，因此与 NMR 不同，无须用脉冲
技巧就能观察自由感应衰减。$\mu$ 子会在磁场——内场或外场——中拉莫尔进动，进动可
通过测量不对称度跟踪。进动频率通过 $\omega=\gamma_{\mu}B$ 直接联系磁场，其中
$\gamma_{\mu}=ge/2m_{\mu}=2\pi\times135.5\ \mathrm{MHz\,T^{-1}}$ 是 $\mu$ 子的旋磁
比（$m_{\mu}$ 为其质量，此处 $g\approx2$）。因此，注入磁有序材料的 $\mu$ 子在内场
中进动，直接给出频率正比于内磁场的振荡信号——在这个意义上，$\mu$ 子是一台微观
磁强计。质子（NMR）、电子（ESR）与 $\mu$ 子（$\mu$SR）的拉莫尔进动频率示于
图 3.24。

\begin{figure}
\centering
\includegraphics[width=0.5\textwidth]{fig3_24.pdf}
\caption{电子、$\mu$ 子与质子的拉莫尔进动频率 $f$（及相应周期 $\tau=1/f$）随外加
磁场 B 的变化。}
\end{figure}

$\mu$ 子磁矩大，对极小的磁场（低至 $\sim10^{-5}$ T）也很敏感，因此适合研究小磁矩
磁性；对磁有序随机或极短程的材料也有价值。由于 $\mu$ 子在样品中均匀停下，每个
信号在实验谱中以正比于其体积分数的强度出现——故对多相或未完全有序的样品尤为
有用。由于（与后面讨论的衍射技术不同）不获得空间信息，单晶样品并非必需（尽管
某些情形有用），而且对某些不便做常规磁中子衍射的材料，$\mu$SR 常能提供磁有序
信息。要从 $\mu$SR 实验提取定量信息，须知道 $\mu$ 子落位（muon-site），这在某些
情形会妨碍对数据的直截了当的解释。（NMR 中核把质子定位得很好，探针位置一目了然。）
通常 $\mu$ 子可占据的间隙位是一小组，且在有利情形只有一个与观察数据一致。这一
技术已在磁性材料中得到广泛应用。

\section*{延伸阅读}

\begin{itemize}[itemsep=0.2em]
\item NMR 的优秀入门可见 P.~J.~Hore，《Nuclear magnetic resonance》，OUP 1995；
B.~Cowan，《Nuclear magnetic resonance and relaxation》，CUP 1997 也极为有用；
NMR 的经典著作是 A.~Abragam，《Principles of nuclear magnetism》，OUP 1961。
\item A.~Abragam 与 B.~Bleaney，《Electron Paramagnetic Resonance of Transition
Ions》，Dover 1986，提供了关于顺磁盐中晶体场与 ESR 实验的大量信息。
\item 晶体场可用所谓的 Stevens 算符处理，见 K.~W.~H.~Stevens，Proc.~Phys.~Soc.~A
\textbf{65}, 209 (1952) 与 M.~T.~Hutchings，Solid State Physics \textbf{16}, 227
(1966)。
\item 穆斯堡尔效应的更多信息见《M\"ossbauer spectroscopy》，D.~P.~E.~Dickson 与
F.~J.~Berry 编，CUP 1986。
\item $\mu$SR 的更多信息见 S.~J.~Blundell，Contemp.~Phys.~\textbf{40}, 175 (1999)。
\end{itemize}

\section*{习题}

\exer{3.1} $\mathrm{Sc^{2+}}$ 离子在 3d 壳层有一个电子。它处于各向异性晶体中，
晶体场对 3d 电子的作用可写为势 $A\hat{l}_z^{2}$。$A>0$ 与 $A<0$ 时 Sc 离子的最低
轨道态各是什么？自旋--轨道耦合 $\lambda\hat{\mathbf{l}}\cdot\hat{\mathbf{s}}$ 远小于
晶体场；把这一项也包括进来后，$A<0$ 与 $A>0$ 时离子的近似基态是什么？讨论沿 $z$
轴以及垂直于 $z$ 轴加小磁场对这些态的影响，并画出磁化率随温度变化的示意图。

\exer{3.2} 把等量的正点电荷放在正八面体的六个顶点上。取笛卡尔坐标系原点位于八面体
中心，证明中心附近势为
\begin{equation*}
V = \frac{q}{4\pi\epsilon_0 a}\left[6 + \frac{35}{4a^{4}}\left(x^{4}+y^{4}+z^{4}-\frac{3}{5}r^{4}\right) + O\left(\frac{r^{6}}{a^{6}}\right)\right],\tag{3.41}
\end{equation*}
$q$ 为每个电荷的大小，$a$ 为原点到每个电荷的距离。

\exer{3.3} 某化合物单位体积含 $n$ 个过渡金属离子，每个离子处于形式为
\begin{equation*}
V = D\left(x^{4}+y^{4}+z^{4}-\frac{3}{5}r^{4}\right)\tag{3.42}
\end{equation*}
的晶体场中，它对简并的 d 能级构成微扰。未微扰波函数可以有多种选取方式；先取
$\hat{L}_z$ 的本征函数并按其本征值 $m_l$ 标记：
\begin{center}
\resizebox{0.88\textwidth}{!}{$\begin{aligned}
&|\pm2\rangle = R(r)\sin^{2}\theta\,\mathrm{e}^{\pm2\mathrm{i}\phi}\\
&|\pm1\rangle = \mp2R(r)\sin\theta\cos\theta\,\mathrm{e}^{\pm\mathrm{i}\phi}\\
&|0\rangle = \sqrt{\tfrac{2}{3}}\,R(r)(3\cos^{2}\theta-1),
\end{aligned}$}\qquad(3.43)
\end{center}
$R(r)$ 是含归一化常数的波函数径向部分。系统的一般态可写为
$|\psi\rangle=\sum_{j=-2}^{2}a_j|j\rangle$，并指定为矢量：
\begin{equation*}
|\psi\rangle = \begin{pmatrix}a_2\\a_1\\a_0\\a_{-1}\\a_{-2}\end{pmatrix}.\tag{3.44}
\end{equation*}
在该基下证明
\begin{equation*}
V = \begin{pmatrix} A & 0 & 0 & 0 & 5A\\ 0 & -4A & 0 & 0 & 0\\ 0 & 0 & 6A & 0 & 0\\ 0 & 0 & 0 & -4A & 0\\ 5A & 0 & 0 & 0 & A \end{pmatrix}\tag{3.45}
\end{equation*}
$A$ 为常数，并证明本征值与本征函数为：
\begin{center}
\begin{tabular}{lccccc}
本征值： & $6A$ & $6A$ & $-4A$ & $-4A$ & $-4A$\\[2pt]
本征函数： & $|0\rangle$ & $\dfrac{|2\rangle+|-2\rangle}{\sqrt{2}}$ & $|1\rangle$ & $|-1\rangle$ & $\dfrac{|2\rangle-|-2\rangle}{\sqrt{2}}$
\end{tabular}\hfill(3.46)
\end{center}
磁场 B 诱导劈裂 $\mu_{\mathrm{B}}Bm_l$，于是 V 变为
\begin{equation*}
V = \begin{pmatrix} A+2\mu_{\mathrm{B}}B & 0 & 0 & 0 & 5A\\ 0 & -4A+\mu_{\mathrm{B}}B & 0 & 0 & 0\\ 0 & 0 & 6A & 0 & 0\\ 0 & 0 & 0 & -4A-\mu_{\mathrm{B}}B & 0\\ 5A & 0 & 0 & 0 & A-2\mu_{\mathrm{B}}B \end{pmatrix}\tag{3.47}
\end{equation*}
计算此时的本征值与本征矢。本征值的图示于图 3.25。

再证明低温下（$k_{\mathrm{B}}T\ll A$），趋于零的磁场中磁化率等于
\begin{equation*}
\chi = \begin{cases} \dfrac{2\mu_0\mu_{\mathrm{B}}^{2}n}{3k_{\mathrm{B}}T} & \text{当 } A>0\\[10pt] \dfrac{2\mu_0\mu_{\mathrm{B}}^{2}n}{5A} & \text{当 } A<0, \end{cases}\tag{3.48}
\end{equation*}
第一种情形不依赖温度且很小，第二种情形等价于 $l=1$ 自由离子的预期值。

\begin{figure}
\centering
\includegraphics[width=0.55\textwidth]{fig3_25.pdf}
\caption{晶体场把 $l=2$ 的离子分裂成一个三重态和一个激发单态；磁场 B 使这些态进一步
分裂。所列本征态指零场情形（场增大时，较高的态逐渐获得更多 $|2\rangle$ 成分，较低
的态逐渐获得更多 $|-2\rangle$ 成分）。}
\end{figure}

\exer{3.4} 某离子核自旋为零，其基态由两个简并能级组成，对应有效自旋 $S=1/2$。
施加磁通密度为 B 的磁场，产生与 B 成线性的能级分裂。观察到该离子的单条电子顺磁
共振线，频率 30 GHz、磁通密度 0.6 T。核自旋非零 I 的同位素产生超精细结构（在自旋哈密顿量
中由 $A\mathbf{I}\cdot\mathbf{S}$ 项描述），由四条近似等距的共振线组成，间隔
$10^{-2}$ T，对称分布在核自旋为零的同位素的谱线两侧。这给出了关于（a）核自旋与
（b）该同位素核磁矩的什么信息？计算参数 A 的数值。

为什么有时必须在低温下进行这些测量？

\exer{3.5} 图 3.26 给出三种顺磁盐每离子磁矩对 $B/T$ 的图。证明式 2.9 在低与高
$B/T$ 处都与数据同形。自由 $\mathrm{Cr^{3+}}$ 与 $\mathrm{Fe^{3+}}$ 离子的 3d 壳层
分别含三与五个电子，$\mathrm{Gd^{3+}}$ 的 4f 壳层含七个电子。用洪特规则计算这些
离子的 $S, L, J$ 与 $g_J$，证明所得数值能很好解释图中 $\mathrm{Gd^{3+}}$ 与
$\mathrm{Fe^{3+}}$ 的数据，但不能解释 $\mathrm{Cr^{3+}}$。从实验数据推断
$\mathrm{Cr^{3+}}$ 的实际 $g$ 因子与基态量子数，并简述数据与你最初计算不符的原因。

\begin{figure}
\centering
\includegraphics[width=0.5\textwidth]{fig3_26.pdf}
\caption{三种顺磁盐的每离子磁矩：$\mathrm{KCr(SO_4)_2\cdot12H_2O}$
（$\mathrm{Cr^{3+}}$）、$\mathrm{NH_4Fe(SO_4)_2\cdot12H_2O}$（$\mathrm{Fe^{3+}}$）
与 $\mathrm{Gd_2(SO_4)_3\cdot8H_2O}$（$\mathrm{Gd^{3+}}$）。取自 W.~Henry，
Phys.~Rev.~\textbf{88}, 559 (1952)。}
\end{figure}

\exer{3.6} 一台 NMR 谱仪工作在 60 MHz。${}^{1}$H、${}^{2}$H、${}^{13}$C 与
${}^{19}$F 核的共振预计出现在多大外加磁场处？（用表 2.3 的数据。）

\exer{3.7} 一台 ESR 谱仪工作在 9 GHz（X 波段）。要观察 DPPH（一种用于校准 ESR
谱仪的有机分子，在 $g=2$ 处给出尖锐信号）中未成对电子的信号，需要多大磁场？

\exer{3.8} 某系统的自旋哈密顿量为
\begin{equation*}
\hat{\mathcal{H}} = g_{\parallel}\mu_{\mathrm{B}}B_zS_z + g_{\perp}\mu_{\mathrm{B}}(B_xS_x+B_yS_y) + DS_z^{2}.\tag{3.49}
\end{equation*}
证明：对该系统中处于三重态（$S=1$）的两个电子，能量本征值 E 可由
\begin{equation*}
E(E-D)^{2} - E(g_{\parallel}\mu_{\mathrm{B}}B\cos\theta)^{2} - (E-D)(g_{\perp}\mu_{\mathrm{B}}B\sin\theta)^{2} = 0.\tag{3.50}
\end{equation*}
得到。求并画出 $\theta=0$ 与 $\theta=\pi/2$ 的解。示意 $\hbar\omega>D$ 时可能的
光子跃迁，进而说明如何通过测量 $\theta=0$ 与 $\theta=\pi/2$ 的 ESR 得到参数
$g_{\parallel}$、$g_{\perp}$ 与 D。

\exer{3.9} 穆斯堡尔效应中 ${}^{57}$Fe 原子发射 14.4 keV 光子。计算自由
${}^{57}$Fe 原子以及被刚性固定在 10 mg 晶体晶格中的 ${}^{57}$Fe 原子的反冲速度。
两种情形下 14.4 keV 光子的多普勒频移各是多少？某穆斯堡尔实验在源与样品相对速度
为 2 mm\,s$^{-1}$ 时得到尖锐谱线。计算相应的频移（Hz 与 cm$^{-1}$）以及能量移动
（eV）。

\exer{3.10} $\mu$ 子平均寿命 2.2 $\mu$s；假定可测比例的 $\mu$ 子能活 20 $\mu$s，
这给可测内场设置了什么下限？一次典型实验测量 $10^{7}$ 个 $\mu$ 子，其中有多少能
活 20 $\mu$s 或更久？同步加速器产生的 $\mu$ 子脉冲宽度为 50 ns。这给样品中可测
内场设置了什么上限？（提示：考虑脉冲前端与后端 $\mu$ 子产生的进动信号的差异。）
回旋加速器可产生连续 $\mu$ 子束，避免这一问题。某磁有序氧化物中被注入 $\mu$ 子
占据的间隙位内场为 0.4 T。预期测到多大的进动频率？

\exer{3.11} 晶体中 $\mathrm{Fe^{3+}}$ 离子（${}^{6}S_{5/2}$）的饱和磁矩预期为
$5\mu_{\mathrm{B}}$（见图 3.26），但由磁化率测量推得的有效磁矩预期为
$5.92\mu_{\mathrm{B}}$（见表 3.1）。为何不同？
